\section{用一把车刀切削正多边形零件}
\markboth{45. 用一把车刀切削正多边形零件}{45. 用一把车刀切削正多边形零件}

在第 44 节中已经证明了，如果让工件和刀盘在同一平面内沿相同方向转动，并且工件转速是刀盘转速的 $m$ 倍，那么装在刀盘上的任一把车刀的刀尖 $M$ 在工件平面内的轨迹，可用下面的方程表示：
\begin{align*}
  x &= l\cos m\theta + h\cos(\theta - m\theta), \\
  y &= l\sin m\theta - h\sin(\theta - m\theta),
\end{align*}
这里 $l$ 是工件中心与刀盘中心的距离，$h$ 是刀尖 $M$ 到刀盘中心的距离，$0<h<l$，并且 $l-h$ 比 $l+h$ 小得多。在第 44 节中选用 $m=\frac{1}{2}$，因而刀尖在工件平面内的轨迹是狭长的椭圆，三个椭圆才能围出一个六边形区域来。所以选用 $m=\frac{1}{2}$ 时必须用三把车刀才能切削出正六边形。

如果我们在夹具上调换传动齿轮，因而改变工件转速与刀盘转速的比值 $m$，其余条件不变，那么刀尖的轨迹就会变成其它的内摆线族或外摆线族曲线，这时往往一条曲线就能围出一个多边形的“核”来，因而只用一把车刀就能切削正多边形工件。

\begin{figure}[htbp]
  \centering
  \begin{minipage}[b]{0.24\textwidth}
    \centering
    \begin{tikzpicture}[scale=0.58]
      \def\m{1/3}
      \def\h{0.76}
      \draw[thick, curveblue, domain=0:1080, samples=260, variable=\t]
        plot ({cos(\m*\t)+\h*cos((1-\m)*\t)},
              {sin(\m*\t)-\h*sin((1-\m)*\t)});
    \end{tikzpicture}
    \caption*{图\ 45.1}
  \end{minipage}
  \hfill
  \begin{minipage}[b]{0.24\textwidth}
    \centering
    \begin{tikzpicture}[scale=0.58]
      \def\m{3/5}
      \def\h{0.76}
      \draw[thick, curveblue, domain=0:1800, samples=340, variable=\t]
        plot ({cos(\m*\t)+\h*cos((1-\m)*\t)},
              {sin(\m*\t)-\h*sin((1-\m)*\t)});
    \end{tikzpicture}
    \caption*{图\ 45.2}
  \end{minipage}
  \hfill
  \begin{minipage}[b]{0.24\textwidth}
    \centering
    \begin{tikzpicture}[scale=0.58]
      \def\m{1/6}
      \def\h{0.68}
      \draw[thick, curveblue, domain=0:2160, samples=360, variable=\t]
        plot ({cos(\m*\t)+\h*cos((1-\m)*\t)},
              {sin(\m*\t)-\h*sin((1-\m)*\t)});
    \end{tikzpicture}
    \caption*{图\ 45.3}
  \end{minipage}
  \hfill
  \begin{minipage}[b]{0.24\textwidth}
    \centering
    \begin{tikzpicture}[scale=0.58]
      \def\m{3/7}
      \def\h{0.78}
      \draw[thick, curveblue, domain=0:2520, samples=420, variable=\t]
        plot ({cos(\m*\t)+\h*cos((1-\m)*\t)},
              {sin(\m*\t)-\h*sin((1-\m)*\t)});
    \end{tikzpicture}
    \caption*{图\ 45.4}
  \end{minipage}
\end{figure}

在图 45.1 $\sim$ 45.4 中分别画出 $m=\frac{1}{3}$、$\frac{3}{5}$、$\frac{1}{6}$ 和 $\frac{3}{7}$ 的情形\footnote{重排校订：原书图 45.1（$m=\frac{1}{3}$）将极半径极大值（顶点）误绘于正左方，与参数方程 $\theta=0$ 时 $(x,y)=(l+h,0)$ 位于正右方的方位相差 $180^\circ$；本版据参数方程真实方位统一重绘图 45.1～45.4。}，这些曲线的中心都有一个多边形“核”，多边形的边数顺次是 $3$、$5$、$6$、$7$。所以，如果让工件转速是刀盘转速的 $\frac{1}{6}$，只要用一把车刀就能切削螺丝帽侧面的正六边形。而如果想切削正五边形的零件，只要改变转速比，使工件转速是刀盘转速的 $\frac{3}{5}$ 就行了。

一般地，如果想用一把车刀切削正 $n$ 边形，应该怎样选择 $m$ 的值呢？

由于采用齿轮传动，转速比是两个正整数的比，因而比值 $m$ 必为有理数。把 $m$ 表示成既约分数 $\frac{q}{p}$，并且把 $m=\frac{q}{p}$ 代入曲线方程，得到
\begin{align*}
  x &= l\cos\frac{q\theta}{p} + h\cos\frac{(p-q)\theta}{p}, \\
  y &= l\sin\frac{q\theta}{p} - h\sin\frac{(p-q)\theta}{p}.
\end{align*}
因为 $q$ 与 $p$ 互质，所以 $p-q$ 也与 $p$ 互质。由此知道 $2p\pi$ 是函数 $\cos\frac{q\theta}{p}$、$\sin\frac{q\theta}{p}$、$\cos\frac{(p-q)\theta}{p}$ 和 $\sin\frac{(p-q)\theta}{p}$ 的周期，因而 $x$ 和 $y$ 也都是周期为 $2p\pi$ 的函数。只要让参数 $\theta$ 在范围 $0\leqslant\theta<2p\pi$ 内变化，就能得到整个曲线。

从曲线方程又得到
\[
  x^2+y^2 = l^2+h^2+2lh\cos\left(\frac{q\theta}{p}+\frac{(p-q)\theta}{p}\right),
\]
即
\begin{align*}
  x^2+y^2 &= l^2+h^2+2lh\cos\theta \\
          &\geqslant l^2+h^2-2lh \\
          &= (l-h)^2.
\end{align*}
由于 $0<h<l$，所以
\[
  \sqrt{x^2+y^2} \geqslant l-h,
\]
等号当且仅当 $\cos\theta=-1$ 时成立。$\sqrt{x^2+y^2}$ 是曲线上的点到坐标原点的距离，使它取得极小值 $l-h$ 的点，是曲线上离它的中心 $O$ 最近的点，即多边形“核”的某一边的中点。因而，多边形的边数 $n$，恰好等于曲线上使 $\sqrt{x^2+y^2}$ 取极小值的点的个数，即在范围 $0\leqslant\theta<2p\pi$ 内满足 $\cos\theta=-1$ 的 $\theta$ 值的个数。满足这些条件的 $\theta$ 值是
\[
  \theta = \pi + 2k\pi,\quad (k=0,\, 1,\, \cdots,\, p-1)
\]
共有 $p$ 个。所以 $n=p$。

这样一来，我们就得到一个很方便的计算法则：要想用一把车刀切削正 $n$ 边形的小零件，只要选取转速比，使工件的转速是刀盘转速的 $\frac{q}{n}$ 倍，这里 $\frac{q}{n}$ 是一个既约分数。

例如，想用一把车刀切削出正六边形轮廓，$m=\frac{q}{n}=\frac{q}{6}$，$q$ 取 $1$ 或 $5$ 时，$\frac{1}{6}$ 和 $\frac{5}{6}$ 都是既约分数，所以只要使工件转速与刀盘转速的比为 $1:6$ 或 $5:6$，都可以达到目的。
